% PASS F — global sign and phasor conventions (Chapter 1)
\paragraph{Global conventions.}
Harmonic fields use $\exp(+\jj\omega t)$ with angular frequency $\omega$ in rad/s, so $\partial_t\rightarrow +\jj\omega$ on phasors.
Outgoing spherical/cylindrical waves carry spatial factors $\exp(-\jj kr)$ or $\exp(-\jj k_z z)$ in the forward-propagation branch.
Maxwell's equations in vacuum SI form read
$\curl\E=-\jj\omega\mu_0\H$,
$\curl\H=\Jimp+\jj\omega\varepsilon_0\E$,
$\divg(\varepsilon_0\E)=\rho$,
$\divg(\mu_0\H)=0$.
The complex Poynting vector is $\mathbf{S}_\mathrm{c}=\tfrac12\E\times\H^*$;
time-averaged power flux is $\langle\mathbf{S}\rangle=\Re(\mathbf{S}_\mathrm{c})$.
Time-averaged momentum density is $\mathbf{g}_{\mathrm{em}}=\varepsilon_0\mu_0\Re(\mathbf{S}_\mathrm{c})$.
The symmetric Maxwell stress tensor $\mathbf{T}$ is given by Eq.~\eqref{eq:ch1.maxwell-stress-components} and Eq.~\eqref{eq:ch1.maxwell-stress}.
Traction on a surface element with outward normal $\hat{\mathbf{n}}$ is $t_i=T_{ij}n_j$, equivalently $\mathbf{t}=\mathbf{T}\cdot\hat{\mathbf{n}}$ \cite{stratton1941}.
The Lorentz force density on matter is $\mathbf{f}_{\mathrm{matter}}=\rho\E+\mathbf{J}\times\B$; on impressed-source support $\mathbf{J}=\Jimp$.
Local momentum balance: $\partial_t\mathbf{g}_{\mathrm{em}}-\divg\mathbf{T}=-\mathbf{f}_{\mathrm{matter}}$;
integral form $\mathbf{F}_{\mathrm{matter}}=\oint \mathbf{T}\cdot\hat{\mathbf{n}}\,dS-(d/dt)\int \mathbf{g}_{\mathrm{em}}\,dV$.
Mechanical torque density on matter is $\boldsymbol{\tau}_{\mathrm{matter}}=\mathbf{r}\times\mathbf{f}_{\mathrm{matter}}$.
The stress moment is $\mathbf{M}=\mathbf{r}\times\mathbf{T}$.
The outward field angular-momentum current is $\mathbf{J}_J=-\mathbf{M}$.
Local angular-momentum balance: $\partial_t\mathbf{j}_{\mathrm{em}}-\divg\mathbf{M}=-\boldsymbol{\tau}_{\mathrm{matter}}$ with $\mathbf{j}_{\mathrm{em}}=\mathbf{r}\times\mathbf{g}_{\mathrm{em}}$, equivalently $\partial_t\mathbf{j}_{\mathrm{em}}+\divg\mathbf{J}_J=-\boldsymbol{\tau}_{\mathrm{matter}}$.
At harmonic steady state on a control volume $V$ with outward $\hat{\mathbf{n}}$,
\begin{equation}
\label{eq:ch1.convention-am-balance}
\frac{d}{dt}\int_V \mathbf{j}_{\mathrm{em}}\,dV
-\oint_{\partial V}\mathbf{M}\cdot\hat{\mathbf{n}}\,dS
=-\int_V \boldsymbol{\tau}_{\mathrm{matter}}\,dV,
\end{equation}
The resultant mechanical torque on matter is
$\mathbf{T}_{\mathrm{matter}}=\int_V \boldsymbol{\tau}_{\mathrm{matter}}\,dV=\oint_{\partial V}\mathbf{r}\times(\mathbf{T}\cdot\hat{\mathbf{n}})\,dS$ at harmonic steady state when storage in $V$ is constant,
and outward field angular-momentum flux is
$\boldsymbol{\Phi}_{J,\mathrm{out}}=-\mathbf{T}_{\mathrm{matter}}$
when the same control volume and origin apply \cite{stratton1941,beth1936}.
The outward $z$-angular-momentum flux density per steradian is
$\Psi_{J_z}^{\mathrm{(out)}}=-r^2[\mathbf{r}\times(\mathbf{T}\cdot\hat{\mathbf{r}})]\cdot\hat{\mathbf{z}}$;
$\dot J_{z,\mathrm{out}}=\int_{4\pi}\Psi_{J_z}^{\mathrm{(out)}}\,d\Omega$ \cite{stratton1941}.
The vector density $\boldsymbol{\Xi}_J=-r^2\,\mathbf{r}\times(\mathbf{T}\cdot\hat{\mathbf{r}})$ carries all three components when needed \cite{beth1936}.
Under $\exp(+\jj\omega t)$, outgoing cylindrical radial waves use $H_m^{(2)}(k_\rho\rho)$ and outgoing spherical radial waves use $h_\ell^{(2)}(kr)=j_\ell(kr)-\jj y_\ell(kr)$; $H_m^{(1)}$ and $h_\ell^{(1)}$ are incoming branches \cite{stratton1941,harrington2001}.
In the radiation zone the outgoing selection implies $\H=\eta_0^{-1}\hat{\mathbf{r}}\times\E$ on $S_R$.
Multipole field pairs assembled from vector Helmholtz solutions must satisfy $\curl\E=-\jj\omega\mu_0\H$ with this branch; when a reference multipole formula is quoted from a source using $\exp(-\jj\omega t)$, the magnetic coefficient may carry an overall sign relative to the electric coefficient even though $h_\ell^{(2)}$ is unchanged \cite{stratton1941,hansen1988}.
Spatial Fourier transforms use
$\tilde{f}(\mathbf{k})=\int f(\mathbf{r})\,e^{\jj\mathbf{k}\cdot\mathbf{r}}\,d^3r$
and
$f(\mathbf{r})=\int \tilde{f}(\mathbf{k})\,e^{-\jj\mathbf{k}\cdot\mathbf{r}}\,d^3k/(2\pi)^3$;
the angular spectrum on $z=0$ uses the two-dimensional specialization of this convention in Eq.~\eqref{eq:ch1.angular-spectrum}.

\noindent\fbox{\parbox{0.97\linewidth}{%
\textbf{Stress and torque convention (Chapter~1).}\;
Traction: $\mathbf{t}=\mathbf{T}\cdot\hat{\mathbf{n}}$ with outward $\hat{\mathbf{n}}$.
$\partial_t\mathbf{g}_{\mathrm{em}}-\divg\mathbf{T}=-\mathbf{f}_{\mathrm{matter}}$.
Stress moment: $\mathbf{M}=\mathbf{r}\times\mathbf{T}$.
Outward field angular-momentum current: $\mathbf{J}_J=-\mathbf{M}$.
Resultant mechanical torque at steady state:
$\mathbf{T}_{\mathrm{matter}}=\oint_{\partial V}\mathbf{r}\times(\mathbf{T}\cdot\hat{\mathbf{n}})\,dS$.
Outward field angular-momentum flux:
$\boldsymbol{\Phi}_{J,\mathrm{out}}=-\mathbf{T}_{\mathrm{matter}}$ when $\dot{\mathbf{J}}_{\mathrm{em}}=0$.
Outward flux density per steradian:
$\Psi_{J_z}^{\mathrm{(out)}}=-r^2[\mathbf{r}\times(\mathbf{T}\cdot\hat{\mathbf{r}})]\cdot\hat{\mathbf{z}}$ with
$\dot J_{z,\mathrm{out}}=\int_{4\pi}\Psi_{J_z}^{\mathrm{(out)}}\,d\Omega$.
}}
